% =============================================================================
%  Vitality-profile methodology — skeleton paper
%  STATUS: UNPROVEN. Methodology paper; algorithmic exposition.
% =============================================================================
\documentclass[11pt]{article}

\usepackage[utf8]{inputenc}
\usepackage[T1]{fontenc}
\usepackage[expansion=false]{microtype}
\usepackage{amsmath, amssymb, amsthm, mathtools}
\usepackage[margin=1in]{geometry}
\usepackage{enumitem}
\usepackage[hidelinks]{hyperref}

\theoremstyle{plain}
\newtheorem{theorem}{Theorem}
\theoremstyle{definition}
\newtheorem{definition}[theorem]{Definition}

\title{Computing the Vitality Profile $(k, \boldsymbol{\sigma})$: \\
       Algorithm, Implementation, and Validation \\
       \large (DRAFT --- skeleton methodology paper)}
\author{Liam McGuire\thanks{Unaffiliated researcher, Seattle, USA. liammcg44@gmail.com}}
\date{\today}

\begin{document}
\maketitle

\begin{abstract}
This paper is a \textbf{skeleton} for a methodology paper presenting
the algorithm for computing the vitality profile $(k, \boldsymbol{\sigma})$
on a \texttt{System} instance, following the algorithmic
specification in \texttt{docs/vitality\_computation.md}. \textbf{No
proof is attempted in the conjecture sense; the methodology itself is
an exposition of an unproven algorithm rather than a mathematical
theorem.} The paper is the operational complement to the framework's
predictive thesis.
\end{abstract}

\section{Goal (UNPROVEN-as-validated)}

Present the vitality-profile algorithm in publishable form, with
sufficient detail for an external researcher to implement it and to
validate it against the regression-test cases in
\texttt{vitality\_computation.md} (hurricane $\to (0, ())$,
converged GA $\to (1, (0))$, active Tierra $\to (2, (1, 1))$, etc.).

\section{Plan of attack (no proof attempted)}

\begin{enumerate}[leftmargin=2em,itemsep=4pt]
  \item \textbf{Algorithm exposition (UNPROVEN).} Present Stage 1
    (detect $k$ recursively), Stage 2 (measure $\sigma_i$ via
    conditional transfer entropy), Stage 3 (composite systems via
    $\boxtimes_\kappa$), Stage 4 (cross-system comparison machinery),
    and Stage 5 (level-transition detection). The current draft of
    \texttt{vitality\_computation.md} is the source.
  \item \textbf{Implementation reference (UNPROVEN).} Describe the
    \texttt{vitality.py} module (pending, see
    \texttt{FRAMING\_PIVOT\_TASKS.md} Part 2 C-5) and how it composes
    existing primitives from \texttt{emergence.py},
    \texttt{entity.py}, \texttt{viability.py}, \texttt{observer.py}.
  \item \textbf{Empirical validation (UNPROVEN).} Run the algorithm
    on the regression-test fixtures. Verify that the verdicts match
    the structural predictions. Where they disagree, identify
    estimator-vs-truth discrepancies vs.\ algorithm errors.
  \item \textbf{Methodology comparison (UNPROVEN).} Compare with
    existing ALife methodology: Phi-ID, Wolfram classes,
    open-endedness metrics. Where the vitality profile gives a
    different verdict, explain the disagreement structurally.
\end{enumerate}

\section{Status and follow-up}

\textbf{Current status: UNPROVEN (as a publishable methodology paper).}
The algorithmic spec exists; the implementation is pending; the
empirical validation requires the implementation to land first. This
paper is the empirical-research-audience-facing exposition once
those pieces are in place.

\bibliographystyle{plain}
\bibliography{references}

\end{document}
